% Chunk: Chapter 8 backmatter---Problems and References.
% Source: Weinberg QFT Vol. I, physical PDF pp. 397--398
%         (printed pp. 374--375).
% Problems: 5. Reference labels: 1, 1a, 2, 3, 4, 5, 6, 7.
% Status: source-reviewed against rendered scans. Uncertainties: none.

\chapterexercisehook{08}

\chapterbackmatter{References}

\begin{enumerate}
\item[1.]\label{ref:8.1} See, e.g., M. B. Green, J. H. Schwarz, and E. Witten,
\emph{Superstring Theory} (Cambridge University Press, Cambridge, 1987):
Section 2.2.

\item[1a.]\label{ref:8.1a} V. Fock, \emph{Z. f. Phys.} \textbf{39}, 226 (1927); H. Weyl,
\emph{Z. Phys.} \textbf{56}, 330 (1929). The term `gauge invariance' derives
from an analogy with earlier speculations about scale invariance by H. Weyl,
in \emph{Raum, Zeit, Materie}, 3rd ed. (Springer-Verlag, Berlin, 1920). Also
see F. London, \emph{Z. f. Phys.} \textbf{42}, 375 (1927). This history has
been reviewed by C. N. Yang, talk at City College (unpublished).

\item[2.]\label{ref:8.2} The use of Coulomb gauge in electrodynamics was strongly advocated
by Schwinger on pretty much the same grounds as here: that we ought not to
introduce photons with helicities other than $\pm1$. See J. Schwinger,
\emph{Phys. Rev.} \textbf{78}, 1439 (1948); \textbf{127}, 324 (1962);
\emph{Nuovo Cimento} \textbf{30}, 278 (1963).

\item[3.]\label{ref:8.3} R. P. Feynman, \emph{Phys. Rev.} \textbf{101}, 769 (1949):
Section 8.

\item[4.]\label{ref:8.4} O. Klein and Y. Nishina, \emph{Z. f. Phys.} \textbf{52}, 853
(1929); Y. Nishina, \emph{ibid.}, 869 (1929); also see I. Tamm,
\emph{Z. f. Phys.} \textbf{62}, 545 (1930).

\item[5.]\label{ref:8.5} See, e.g., S. Weinberg, \emph{Gravitation and Cosmology} (Wiley,
New York, 1972): Section 4.11.

\item[6.]\label{ref:8.6} For a readable general introduction to the geometry and topology
of $p$-forms, see H. Flanders, \emph{Differential Forms} (Academic Press,
New York, 1963).

\item[7.]\label{ref:8.7} T. West, \emph{Comput. Phys. Commun.} \textbf{77}, 286 (1993).
\end{enumerate}
